\ChapterImagePrelim[cap:glosario]{Glosario}{./images/fondo.png}\label{cap:glosario}
\mbox{}\\
En este apartado se encuentran términos clave y conceptos relevantes utilizados a lo largo de este proyecto.

\section*{A}
\begin{description}
	\item[API Gateway:] Componente de la capa de aplicación en la propuesta~\NDITI que centraliza el acceso a servicios y gestiona el enrutamiento de solicitudes.

	\item[Arquitectura multicapa:] Enfoque de la propuesta~\NDITI que organiza el sistema en capas (física, redes, tecnológica, aplicación y negocio) para facilitar su modelado.
\end{description}

\section*{C}
\begin{description}
	\item[Capa de aplicación:] Nivel donde se modela la lógica funcional del software, sus servicios e interacciones en la propuesta~\NDITI.

	\item[Capa de negocio:] Nivel más alto de abstracción que representa procesos, actores y objetivos organizacionales en la propuesta~\NDITI.

	\item[Capa física:] Nivel base del modelo que representa la infraestructura tangible como hardware, redes físicas y recursos materiales.

	\item[Capa tecnológica:] Nivel intermedio que representa entornos de ejecución, virtualización y servicios que soportan aplicaciones.
\end{description}

\section*{D}
\begin{description}
	\item[DAR:] Técnica del modelo CMMI utilizada para evaluar alternativas de diseño y apoyar la toma de decisiones en la construcción de la notación~\NDITI.
\end{description}

\section*{E}
\begin{description}
	\item[Estudio de mapeo sistemático (\SMS):] Método de revisión estructurada de literatura utilizado para identificar y analizar trabajos relacionados con notaciones de modelado.
\end{description}

\section*{I}
\begin{description}
	\item[Infraestructura de TI:] Conjunto de recursos físicos, lógicos y de red que soportan el funcionamiento de sistemas de información.

	\item[Interoperabilidad:] Propiedad que permite a sistemas o modelos distintos intercambiar información de forma coherente; en el contexto de~\NDITI{}, se refiere a la integración entre capas.
\end{description}

\section*{N}
\begin{description}
	\item[NDITI:] Notación propuesta en este trabajo para el diseño de Infraestructura de TI, basada en ArchiMate con extensiones de UML y BPMN.
\end{description}

\section*{P}
\begin{description}
	\item[Paralelización:] Mecanismo que permite la ejecución simultánea de actividades dentro de un proceso en la notación~\NDITI.

	\item[Proceso de negocio:] Secuencia de actividades que representa la lógica operativa de una organización dentro de la capa de negocio.
\end{description}

\section*{R}
\begin{description}
	\item[Redes y comunicaciones:] Capa de la propuesta~\NDITI que modela dispositivos de interconexión y transmisión de datos.

	\item[Relación:] Conector que define la interacción o dependencia entre elementos dentro del modelo~\NDITI.
\end{description}

\section*{S}
\begin{description}
	\item[Servicio:] Funcionalidad expuesta por componentes de la arquitectura en la propuesta~\NDITI.

	\item[\SMS:] Estudio de mapeo sistemático utilizado para estructurar el análisis de literatura base del proyecto.
\end{description}

\section*{U}
\begin{description}
	\item[UML:] Lenguaje de modelado utilizado como referencia en la propuesta~\NDITI para complementar la representación estructural y conductual de sistemas.
\end{description}